\cotefile{156-5}
\scanpage{10}\mspage{9}
en question sont dans $\Phi$, i.e. que $\Phi$ est \unsure{fermé|fermée} dans $\Phi\cup\Psi$ (et itou pour le cas $A\in\Psi$, avec $\Psi$ fermé dans $\Phi\cup\Psi$). Donc il \unsure{vient|reste}

\[
\begin{aligned}
2^{\circ})\quad &\deleted{\unsure{\Psi}}\Phi\ \text{fermé dans }\Phi\cup\Psi:\quad A\in\Phi,\ B\in\Psi,\\
&B\subset\deleted{\unsure{B}}A\ \Longrightarrow\ B\in\deleted{B}\Phi.
\end{aligned}
\]
Mais on aura $\deleted{\unsure{\lvert B\rvert}}$ $B\subset\lvert\Phi\rvert\cap\lvert\Psi\rvert$, donc $B\in\Phi\Rightarrow B\deleted{\unsure{\subset}}\in\deleted{\unsure{\Phi}}\Psi$ OK.

\uline{\unsure{Définition|Definition}} On dit alors que $\Phi$ et $\Psi$ \uline{compatibles}

\uline{Corollaire 1} Quand ces conditions sont satisfaites, $\lvert\deleted{A}\Phi\rvert\cap\lvert\Psi\rvert$ est une partie de $\lvert\Phi\rvert$ qui est $\Phi$-saturée, et une partie de $\Psi$ qui est $\Psi$-saturée, et les deux relations d'\unsure{équivalence|équilence} induites par $\lvert\deleted{\unsure{B}}\Phi\rvert$ et par $\lvert\Psi\rvert$ sont les $=\mathrm{s}$.

Mais attention, cette dernière condition ne suffit nullement à \unsure{impliquer} que $\Phi$ et $\Psi$ compatibles. Exemple : $\Phi$ et la famille \inserted{$\Phi^{\circ}$} des $A^{\circ}$ ont $=$ support, et définissent la $=$ relation d'équiv. sur $\lvert\Phi\rvert$, mais elles ne sont compatibles que si elles sont \uline{égales} \deleted{(car} i.e. $\Phi$ formé d'\unsure{ens.} deux à deux disjoints. En effet \deleted{(car la relation d'ordre sur $\Phi$ et sur $\Phi^{\circ}$ n'est la $=$ que si}

\uline{Corollaire 2} Supposons $\Phi$ et $\Psi$ compatibles. Alors $\lvert\Phi\rvert\subset\lvert\Psi\rvert$ \unsure{sss|ssi} $\Phi\subset\Psi$, $\lvert\Phi\rvert=\lvert\Psi\rvert$ sss $\Phi=\Psi$.
\scanpage{11}\mspage{10}
En \deleted{\unsure{l}}effet, c'est comme pour les sous-familles admissibles d'une famille (ici $\Phi\cup\deleted{\unsure{\Phi}}\Psi$) -- cf. Scholie page 7.

\margin{\uline{Cor} 3 Si $\Phi,\Psi$ comp., alors pour $\Phi',\deleted{\unsure{\Psi}}\leq\Phi,\ \Psi'\leq\Psi$ (parties fermées) $\Phi'$ et $\Psi'$, \unsure{st} admissibles, et $\Phi'\cup\Psi'$ est fermé (ss-famille adm.) dans $\Phi\cup\Psi$.}

% slanted
\margin{On appellera \unsure{ainsi|aussi} les éléments de $\unsure{\mathcal F|\mathscr F}$ les \uline{``figures \unsure{ensemblistes}''} de $L$}

Soit $\mathcal F$ l'ens des \deleted{familles} \unsure{des|les} parties \uline{admissibles} $\Phi$ de $\unsure{\mathcal P|\mathfrak P}(L)$. On ordonne $\mathcal F$ par
\[
\begin{aligned}
\Phi\leq\Psi\quad\overset{\mathrm{def}}{\Longleftrightarrow}\quad &\Phi\ \text{est une ss-}\\
&\text{famille admissible}\\
&\text{de }\Psi.
\end{aligned}
\]

Alors on a

\uline{C 1} $\mathcal F$ admet des Sup majorés, car si $\Phi_i\in\mathcal F,\ \Phi_i\leq\deleted{\Phi}\Psi$ des ss-ens. fermés d'un $=$ $\deleted{\Phi}\Psi$, alors $\bigcup_i\Phi_i$ est un sous-ens. fermé de $\Psi$, qui est \unsure{le|la} la sup. C'est immédiat par \unsure{le sorite|les sorites} sur les sous-familles \unsure{équiv.} (\unsure{de ses|i.e.|les|deux} parties fermées) d'un $\Psi$

% slanted
\margin{La Sup pour la relation $\leq$ (qui est ici $\subset$) est $\bigcup_i$}

\deleted{\unsure{Oui|Or|O.K.}, un peu \unsure{vif|cf.}, il faut \unsure{passer|préciser} par des Sup finis, en notant que $\emptyset\subset\mathcal P(L)$ est admissible (! cas du Sup \inserted{de la \unsure{fam.}} vide !) -- la seule famille admissible $\underset{\in\mathcal F}{\Phi}$ telle que $\lvert\Phi\rvert=\emptyset$ et que si $\Phi,\Psi,\Sigma$ sont tels que}

% slanted
\margin{\unsure{explicitons|expliciter} \unsure{cf.} \uuline{Prop} ou Cor}

% This line has a horizontal cancellation in addition to the diagonal strokes across the draft.
\deleted{$\Phi$ et $\Psi$) est admissible, alors $\Sigma$ \unsure{ai|est}}

\deleted{($\Phi,\Psi$) compatibles, alors $\Sigma$ est compatible avec $\Phi\unsure{\cup|et}\Psi$ \unsure{sss|ssi} il l'est avec $\Phi$ et avec $\Psi$}
\scanpage{12}\mspage{11}
Les éléments irréductibles de $\unsure{\mathcal F|\mathscr F}$ sont les $\Phi\in\mathcal F$ qui ont \deleted{\unsure{une}} un plus grand élément : car un \unsure{ens} fermé est irréductible (au sens fort) \unsure{sss|ssi} il a un plus grand élément.

% slanted
\margin{Appelons-les les ``\unsure{multistrates} ensemblistes''}

Donc

\uline{C 2} Tout $\Phi\in\deleted{\unsure{\Phi}}\mathcal F$ est la sup des objets irréductibles de $\mathcal F$, à savoir des $\Phi_{X}=\{Y\in\Phi\mid Y\unsure{\subseteq|\leq}X\}$, pour $Y\in\Phi$

\uline{C 3} Soient $\deleted{\Phi}F,\deleted{\Psi}G\in\mathcal F$. Pour que $F$ et $G$ soient compatibles (i.e $F,G$ \unsure{majoré|majorés}) il f. et s. que $\deleted{X\in F,\unsure{\Rightarrow|\to}}$ $\forall\,\deleted{\Phi}\unsure{X|{\Phi}'}\leq F$ et $\deleted{\Psi}Y\leq G$, $\deleted{\Phi}X$ et $\deleted{\Psi}Y$ \deleted{\unsure{élé|él.}} multistrates, on ait $\deleted{\Phi}\deleted{\Phi}X$ et $\deleted{\Psi}Y$ compatibles.

En d'autres termes, que pour $X\in F,\ Y\in G$, $F_X$ et $G_Y$ soient compatibles :

\uline{\unsure{Remords|Remonter|Remontons} à prop. page \deleted{\unsure{9}}7} : Cond. équiv à a) b) c)

d) \deleted{\unsure{lorsque|toujours|comp.}} $\forall\,A\in\Phi,\ B\in\Psi$, $A\cap B$ est fermé dans $A$ (pour $\Phi$ i.e pour $\Phi_A$) et en $B$ (pour $\Psi$ i.e. pour $\Psi_B$), et si $\deleted{\unsure{A}\ C\in\Phi\cup\Psi}$ $C$ est une partie de $A\cap B$, alors $C\in\Phi$ ($\Leftrightarrow C\in\Phi_A$) sss $C\in\Psi$ ($\Leftrightarrow C\in\Psi_B$) i.e $\Phi_{A\cap B}=\Psi_{A\cap B}$
\scanpage{13}\mspage{12}
\deleted{ce qui}

Composantes connexes de $M=$ ens des multistrates. Les éléments minimaux de $M$ sont ceux qui sont réduits à un seul ensemble $A$, \deleted{(qui peut être} $A\ne\emptyset$ (NB $\Phi=\{\emptyset\}$ n'est pas une famille admissible !\deleted{)} Si $\Phi$ est admissible, $\emptyset\notin\deleted{\unsure{F}}\Phi$, \deleted{)} puisque $A\in\Phi\Rightarrow A^{\circ}\ne\emptyset$ a fortiori $A\ne\emptyset$)
\[
\Phi\deleted{\unsure{F_{A}|\Phi_{A}}}=\{A\}
\]

\deleted{(\unsure{NB})} \deleted{(NB $F\in\unsure{F|M}\Rightarrow F\ne\emptyset$ !)}

\unsure{Si|Soit} $X\deleted{\unsure{\Phi|F}}\in M$, \deleted{et $A\in\deleted{\Phi}F$, alors \unsure{pour}} \inserted{$\deleted{\unsure{F}}X\subset\unsure{\mathcal P|\mathfrak P}(L)$}, et si $\deleted{F}X$ admet un \deleted{plus petit} élément minimal $A$ (p.ex. $\deleted{F}X$ fini, \inserted{on suppose $X\ne\emptyset$}) alors,
\[
\deleted{F}X_A=\{A\}\leq F
\]
\deleted{d'autre part,} donc $\deleted{F_A}$ \deleted{et} $\deleted{F}X$ et $\{A\}$ dans la $=$ comp. connexe. D'autre part, \deleted{\unsure{$\{A\}$ est}} si $A\ne L$, et si $x\in L\setminus A$, $\{A\}$ et $\{\{x\}\}$ sont dans la $=$ composante connexe, car ils sont majorés par le multistrate
\[
\bigl\{\{A\},\{\{x\}\},A\cup\{\{x\}\}\bigr\}
\]
\deleted{donc} Et pour la $=$ raison, \deleted{si} \deleted{\unsure{Card|tout}} \unsure{on} voit que \inserted{$\{\{x\}\}$ pour tous les} \deleted{\unsure{tous les}} \unsure{pts} de $L$ sont dans la $=$ comp connexe \unsure{de $L$|de $\mathscr L$}. En résumé les $X\in M$ qui ont un plus petit élément \inserted{$A$} sont dans une $=$ composante connexe, du moment que $A\ne L$ i.e. du moment que $X$ ne se réduise à $\{A\}=\{L\}$. Mais celui-ci aussi, si \deleted{\unsure{Card}} $L\deleted{\unsure{\geq}}\ne\deleted{\unsure{2|0}}\emptyset$, est dans la
\scanpage{14}\mspage{13}
$=$ comp. connexe que les $\{\{x\}\}$, comme on voit à l'aide de \unsure{la|le} multistrate
\[
\bigl\{\{\{x\}\},\{L\setminus\{x\}\}\bigr\}
\]

Si $L$ \deleted{de card. $1$,} $=\emptyset$, \deleted{$\unsure{\mathcal P|\mathfrak P}(L)$ est de} \deleted{\unsure{cardi}} alors $M=\emptyset$, et $M$ est connexe.

Si $\operatorname{Card}L=1$, $\mathcal P(L)=\{\emptyset,L\}$, et les $\Phi\unsure{\subset|\in}\mathcal P(L)$ \deleted{qui sont \unsure{des} sont que} \inserted{telles} \deleted{\unsure{$\{\Phi\}$} qui ne contiennent pas $\emptyset$} \unsure{que|pour} $\emptyset\notin\Phi$ sont $\emptyset$ et $\{L\}$. Donc $M=\{\{L\}\}$ a un seul élément, et est connexe.

Donc

\uline{Prop} Si $\operatorname{card}L\leq1$, $M$ est connexe (il est vide si $\deleted{\unsure{\mathrm{Card}}}L=\emptyset$, il est de card $1$ donc $0$-connexe si $\operatorname{card}L=1$). Si $\operatorname{card}L\geq1$, alors les $\deleted{\unsure{Y|F}}X\in M$ qui ont un élément minimal sont tous dans la $=$ comp. connexe. Donc si on se borne aux \deleted{familles} figures ensemblistes \uline{finies} (ce qui n'offre pas d'inconvénient à mes yeux) $M$ est connexe. Donc l'axiome C4 \deleted{est ici de} relatif à la décomposition de $M$ en comp. connexes (que $X,Y$ appartenant à des composantes différentes sont \unsure{compatible|compatibles}) est automatique.
\scanpage{15}\mspage{14}
Si on ne se borne aux familles finies, \unsure{Con.|on a|ben.} encore connexité. En effet, si $X\in M$ est $\ne\{L\}$, soit $A\in X$ \unsure{tels|tel|telle} que $A\ne L$, Alors $X$ et $X_A$ dans $=$ comp \inserted{connexe et}
\[
\lvert X_A\rvert=A
\]
donc si $\deleted{\{x\}}x\in L\setminus A$, on trouve une multistrate en \unsure{prenant|posant}
\[
Z=\bigl\{B\unsure{,|;}B\cup\{x\},\{x\}\mid B\in X_A\bigr\}
\]
\margin{\unsure{,}}

\uline{NB} Construction qu'on a \unsure{refait|refaite} ici plusieurs fois : Soient $\deleted{X}F\in\deleted{M}\unsure{\mathcal F|\mathscr F}$, $\deleted{Y}G\in\deleted{\unsure{\mathcal G|M}}\unsure{\mathcal F|\mathcal G}$, $\lvert\deleted{\unsure{X}}F\rvert\cap\lvert G\rvert=\emptyset$, posons $F\boxtimes G=\{A\cup B\mid A\in\deleted{\unsure{X}}F,\ B\in G,\ \inserted{\text{$A$ et $B$ pas tous deux $\emptyset$}}\}$, alors $(A\cup B)^{\circ}=A^{\circ}\cup B^{\circ}$, $\partial(A\cup B)=(\partial A\cup B)\cup(A\cup\partial B)$ -- si $F$ et $G$ \deleted{sont} \deleted{des} multistrates i.e. \deleted{\unsure{admet|admettent|alors}} \deleted{\unsure{pour|par|pas}} plus grand éléments \deleted{\unsure{et que si}} $A_0,B_0$ alors $A_0\cup B_0$ est pl grd élément de $F\boxtimes G$, qui est une multistrate.

\deleted{Celle-ci contient \unsure{en|un} \unsure{l}} $F\boxtimes G$ contient $F\unsure{\amalg|\sqcup}G$ comme sous-figure. Donc, si $F$ et $G$ sont des multistrates ``\unsure{disjointes}'' elles sont dans la $=$ comp. connexe car \unsure{majorées|disjointes|dominées} par $F\boxtimes G$. Donc en tous cas, \margin{\unsure{,|)}} $M$ connexe ($=$ sans \deleted{\unsure{la|les}} \unsure{limiter|limites} aux \deleted{familles} figures ensemblistes finies).
\scanpage{16}\mspage{15}
Supports et cosupports.

$F$ et $G$ \unsure{de|dans} $\unsure{\mathcal F|\mathscr F}$ sont \uline{disjoints} \inserted{(au sens de $\leq$)} \unsure{sss|ssi} $\lvert F\rvert\cap\lvert G\rvert=\emptyset$. Donc les $F$ disjoints d'une famille de figures $G_i$, sont celles qui \deleted{sont} satisfont $\lvert F\rvert\subset L'$, où $L'=L\setminus\bigcup_i\lvert G_i\rvert$. Donc, on voit que les supports correspondent aux parties quelconques $L'$ de $L$, la \unsure{relation|notion} de complémentarité est la \unsure{notion|relation} ensembliste.

\deleted{Le $\mathcal F(L'$} L'ens des figures \inserted{de $L'$} de support $\subset L'$ \deleted{est} s'identifie aux figures de $L'$, i.e. on trouve les $\mathcal F(L')\subset\mathcal F(L)$ comme sous-ens. saturés (au sens de $\leq$ et de la disjonction \ldots).
% Paper watermark, sideways, partial impressions: RONG / EXTRA STR…
\scanpage{17}\mspage{16}
On récupère l'\unsure{ens} $\unsure{\mathcal P|\mathfrak P}(L)$ en termes de $L_{\leq}$, ou de $(M_{\leq},\text{compatibilité})$,

\margin{\uline{NB} $\unsure{\mathcal F|\mathscr F}$ est formé des \uline{toutes} les parties fermées de $\deleted{\unsure{\mathcal F}}M$ dont deux \unsure{élémts} quelc. sont comp.}

et on retrouve la relation d'ordre sur $\mathcal P(L)$ en termes de la relation d'ordre sur les supports, et \unsure{itou|item|idem} pour \unsure{les passages|le passage} au complémentaire. Mais cela ne redonne pas directement $L$, mais $\widehat L$, le compactifié de Stone ($L$ se récupère \inserted{cependant} comme l'\unsure{ens} des points \uline{isolés} de $\widehat L$). Ceci suggère de \unsure{refaire|faire} \unsure{théorie|suivre} en termes d'un espace compact \deleted{\unsure{tot}} discontinu $X$, et \deleted{\unsure{de}} prendre $\mathcal F$ familles admissibles \unsure{dans|d'ens} $X$ fermés et \unsure{même} à la fois \uline{ouverts} et \uline{fermés}.

Si on ne \unsure{permet|prend} que les familles finies, alors les supports \deleted{\unsure{et}} \inserted{et les \unsure{cosupports|supports} des telles \deleted{familles} figures} s'identifient aux parties à la fois ouv. et f. de $X$. Mais
% slanted
\margin{\unsure{ou d'une|ou d'ens} \unsure{fam.|famille|finie} de telles figures}
si on permet des figures infinies ou des ens de figures infinis on les \unsure{change} en retrouvant \deleted{\unsure{à nouveau}} \inserted{les \unsure{intérieurs}} des parties \inserted{\uline{fermées}} quelconques de $\unsure{X|L}$, \unsure{ou|pour} les \deleted{paires} couples support-cosupport s'identifient aux \deleted{paires} \inserted{couples} $(U,V)$ d'ouverts de $X$,
\[
\begin{aligned}
& U=\operatorname{int}(\overbrace{\unsure{\mathrm f|\mathrm F}U}^{\inserted{\overline U}})=\operatorname{ext}V,\quad
 V=\operatorname{int}(\mathrm f\,\unsure{U|V})=\operatorname{ext}U,\\
& \overset{\inserted{\unsure{-\ \smile\ \mathrm m|-\ \smile\ \mathrm{en}}}}{\overline U\cap\overline V} = X\setminus\{U\cup\unsure{V|X}),
\end{aligned}
\]
i.e. les \unsure{ouvert|ouverts} $U$ tels que $\unsure{\overline U|U}=\operatorname{int}(\overline U)$
